% GENERATED FILE. Edit canonical lessons, then run npm run generate:books.
% canonical-source-hash: fnv1a64:b506df6a4e60ed91
% canonical-lessons: PA-C10-paani, PA-C10-chaa, PA-C10-dudh, PA-C10-roti

\chapter{Water, Tea, Milk, and Bread}
\label{ch:pa-request-food-drink}
\hlchaptermodality{10}

\begin{chapteropening}
\textbf{By the end of this chapter:} \emph{I can politely ask for water, tea, milk, or bread in Punjabi, name which of the four is a loanword and which is inherited with an unknown root, and hear the difference between a tone that stays level and one that never formed at all.}

The chapter closes on \pa{ਰੋਟੀ}, the request that breaks the drink pattern: the reader produces all four requests with \pa{ਕਿਰਪਾ} \pa{ਕਰਕੇ}, sorts \pa{ਪਾਣੀ} and \pa{ਦੁੱਧ} (inherited, known root) against \pa{ਚਾਹ} (loaned from Chinese by the overland route, carrying Punjabi's first noun-level tone) against \pa{ਰੋਟੀ} (inherited, root unknown), and names the first Punjabi chapter to realize polite requesting at all.
\end{chapteropening}

\section[pāṇī]{\pa{ਪਾਣੀ} — the first thing you'll actually ask for}

\label{lesson:PA-C10-paani}



\begin{quote}
\emph{Pasand} let you say what pleases you; \emph{madad karnā} let you ask for help. This chapter asks for something smaller and far more frequent: a drink.
\end{quote}

\subsection*{You'll want to know first — one word}
\begin{quote}
\textbf{\pa{ਪਾਣੀ}} — \emph{pāṇī} — \textbf{water}
\end{quote}

\subsection*{The letters in this word}
\textbf{\pa{ਪ}} \emph{p}, \textbf{\pa{ਾ}} long \emph{ā}, then \textbf{\pa{ਣ}} — the curled-back \emph{ṇ} you already read in \textbf{\pa{ਜਾਣਨਾ}} — and \textbf{\pa{ੀ}} long \emph{ī}. Nothing here is new; every letter is doing work you have already made it do.

\begin{grammarlens}[title={Grammar Lens: asking politely}]
\begin{quote}
\textbf{\pa{ਪਾਣੀ}, \pa{ਕਿਰਪਾ} \pa{ਕਰਕੇ।}} — \emph{pāṇī, kirpā karke.} — \textbf{Water, please.}
\end{quote}

\textbf{\pa{ਕਿਰਪਾ} \pa{ਕਰਕੇ}} \emph{kirpā karke} is built from \textbf{\pa{ਕਿਰਪਾ}}, \textquotedblleft{}kindness, grace,\textquotedblright{} plus \textbf{\pa{ਕਰਕੇ}}, \textquotedblleft{}having done\textquotedblright{} — the same \textbf{\pa{ਕਰ}-} you already own from \textbf{\pa{ਮਦਦ} \pa{ਕਰਨਾ}}, \textquotedblleft{}to do a favour of help,\textquotedblright{} now in a \textquotedblleft{}having done\textquotedblright{} shape instead of the infinitive \textbf{-\pa{ਣਾ}} one. Literally: \textquotedblleft{}having done a kindness.\textquotedblright{} Name the item, add this, and the request is complete — no separate verb for \textquotedblleft{}give\textquotedblright{} needed.
\end{grammarlens}

\begin{cousinweb}
\textbf{pāṇī} comes from Sanskrit \textbf{pānīya}, \textquotedblleft{}that which is to be drunk,\textquotedblright{} built on \textbf{pā-}, \textquotedblleft{}to drink.\textquotedblright{} That root is Indo-European *\emph{peh\textsubscript{3}(i)-}, and English holds two disguised pieces of it: \textbf{potion} and \textbf{poison}, both through Latin \emph{potare}, \textquotedblleft{}to drink\textquotedblright{} — a poison is, etymologically, just a drink gone wrong. Greek made the same root into \emph{sympotēs}, \textquotedblleft{}one who drinks together,\textquotedblright{} which English keeps as \textbf{symposium}.

\emph{Pāṇī} is a noun, and it has no room for a liker or a drinker built into it — the same shape of gap \textbf{\pa{ਪਸੰਦ}} left for you to fill with \textbf{\pa{ਮੈਨੂੰ}}. Ask \textbf{\pa{ਮੈਨੂੰ} \pa{ਪਾਣੀ} \pa{ਪਸੰਦ} \pa{ਹੈ}}, \textquotedblleft{}I like water,\textquotedblright{} and you are running last chapter's grammar on this chapter's first noun.
\end{cousinweb}

\subsection*{Your turn}
\begin{itemize}
\raggedright
  \item \emph{Say it:} \textbf{pāṇī} — water
  \item \emph{Say it:} \textbf{pāṇī, kirpā karke.} — water, please
  \item \emph{Match:} \textbf{karke} against \textbf{madad karnā}'s own \textbf{kar-} root
  \item \emph{Trace it:} \emph{pā-} \textquotedblleft{}to drink\textquotedblright{} $\to$ \textbf{pāṇī}, and $\to$ \textbf{potion}, \textbf{poison}, \textbf{symposium}
  \item \emph{Run through:} \textbf{mainū̃ pāṇī pasand hai}, then \textbf{pāṇī, kirpā karke} — liking and requesting, side by side
\end{itemize}

\begin{tcolorbox}[breakable,colback=teal!4,colframe=teal!35!black,title={Before you move on}]
Ask for water politely. (\emph{Pāṇī, kirpā karke.}) What does \emph{karke} share with \emph{madad karnā}? (The same \textbf{kar-} root, \textquotedblleft{}to do.\textquotedblright{}) Which two English words for a drink, and one for a whole gathering of them, share \emph{pāṇī}'s root? (\textbf{Potion}, \textbf{poison}, \textbf{symposium}.)

Sources: \href{https://en.wiktionary.org/wiki/\%E0\%A8\%AA\%E0\%A8\%BE\%E0\%A8\%A3\%E0\%A9\%80}{Wiktionary: \pa{ਪਾਣੀ}}; \href{https://en.wiktionary.org/wiki/\%E0\%A8\%95\%E0\%A8\%BF\%E0\%A8\%B0\%E0\%A8\%AA\%E0\%A8\%BE}{Wiktionary: \pa{ਕਿਰਪਾ}}.
\end{tcolorbox}
\section[cāh]{\pa{ਚਾਹ} — the \pa{ਹ} that is not a consonant}

\label{lesson:PA-C10-chaa}



\begin{quote}
\emph{Pāṇī, kirpā karke} — water, please. The next drink uses the exact same pattern, and its spelling hides a tone you have not met in a noun yet.
\end{quote}

\subsection*{You'll want to know first — one word}
\begin{quote}
\textbf{\pa{ਚਾਹ}} — \emph{cāh} — \textbf{tea}
\end{quote}

\begin{quote}
\textbf{\pa{ਚਾਹ}, \pa{ਕਿਰਪਾ} \pa{ਕਰਕੇ।}} — \emph{cāh, kirpā karke.} — \textbf{Tea, please.}
\end{quote}

\begin{sounds}
\textbf{\pa{ਚਾਹ}} is spelled with a final \textbf{\pa{ਹ}}, but that \textbf{\pa{ਹ}} is not said as a consonant. Chapter 8 showed a \emph{jh} losing its breath and leaving a \textbf{high, falling} pitch on the vowel before it — falling, because more of the word still followed. Here nothing follows: \textbf{\pa{ਚਾਹ}} is one syllable, so the pitch has nowhere to fall \emph{to}. It simply stays \textbf{high}: \emph{c\'{\={a}}h}, level and high from the vowel to the end of the word.

Gurmukhi has no separate tone mark. A trailing \textbf{\pa{ਹ}} that is not pronounced is, in practice, how the script writes this tone — the same way \textbf{\pa{ਗ਼}}'s dot writes a borrowed sound. Say \textbf{\pa{ਖਾਣਾ}} \emph{khāṇā} again, flat, then \textbf{\pa{ਚਾਹ}} \emph{c\'{\={a}}h}, high: two words that both end their vowel clean, one levelled and one not.
\end{sounds}

\begin{cousinweb}
\textbf{\pa{ਚਾਹ}} is a loanword — unlike \textbf{pāṇī}, which Punjabi has owned since Sanskrit. It came through Classical Persian \textbf{čā}, from Mandarin Chinese \textbf{chá}: the \textbf{overland} route, out through Central Asia and Persia, the same road that gave Hindi-Urdu \emph{chāy}, Russian \emph{chai}, and Turkish \emph{çay}. Chinese merchants on the coast spoke a different dialect, Hokkien, where the same word is \emph{tê} — carried by sea into Dutch and then English \textbf{tea}, French \emph{thé}, German \emph{Tee}. One Chinese word, two pronunciations, sorted by whether a language met it on a caravan route or a ship's deck.
\end{cousinweb}

\subsection*{Your turn}
\begin{itemize}
\raggedright
  \item \emph{Say it:} \textbf{cāh} — tea — high and level, not falling
  \item \emph{Run through:} both requests — \textbf{pāṇī, kirpā karke}; \textbf{cāh, kirpā karke}
  \item \emph{Contrast:} \textbf{khāṇā}, flat; \textbf{cāh}, high
  \item \emph{Trace it:} Chinese \emph{chá} $\to$ Persian \emph{čā} $\to$ \textbf{\pa{ਚਾਹ}}, the overland road
  \item \emph{Recall:} what \textbf{\pa{ਮੈਨੂੰ} ... \pa{ਪਸੰਦ} \pa{ਹੈ}} does to the thing you like
\end{itemize}

\begin{tcolorbox}[breakable,colback=teal!4,colframe=teal!35!black,title={Before you move on}]
Did \emph{cāh} come by land or by sea? (\textbf{Land} — through Persian, from Mandarin.) Why is \emph{cāh}'s tone level instead of falling, unlike \emph{samajhṇā}'s? (\textbf{Nothing follows the vowel} — the word ends right there.) What does \emph{cāh}'s final \textbf{\pa{ਹ}} actually mark, if not the sound \emph{h}? (\textbf{The tone itself.})

Sources: \href{https://en.wiktionary.org/wiki/\%E0\%A8\%9A\%E0\%A8\%BE\%E0\%A8\%B9}{Wiktionary: \pa{ਚਾਹ}}; \href{https://en.wiktionary.org/wiki/chai}{Wiktionary: chai}.
\end{tcolorbox}
\section[dudh]{\pa{ਦੁੱਧ} — a third request, and a tone that stays away}

\label{lesson:PA-C10-dudh}



\begin{quote}
\textbf{\pa{ਪਾਣੀ}} carried no tone; \textbf{\pa{ਚਾਹ}} carried a level high one. This word settles which is the ordinary case.
\end{quote}

\subsection*{You'll want to know first — one word}
\begin{quote}
\textbf{\pa{ਦੁੱਧ}} — \emph{dudh} — \textbf{milk}
\end{quote}

\begin{quote}
\textbf{\pa{ਦੁੱਧ}, \pa{ਕਿਰਪਾ} \pa{ਕਰਕੇ।}} — \emph{dudh, kirpā karke.} — \textbf{Milk, please.}
\end{quote}

\subsection*{The letters in this word}
\textbf{\pa{ਦ}} \emph{d}, \textbf{\pa{ੁ}} short \emph{u}, then the \textbf{\pa{ੱ}} \emph{addak} doubling the consonant that follows — the same mark you already read in \textbf{\pa{ਰੱਬ}} — and \textbf{\pa{ਧ}} \emph{dh}, the breathy \emph{d} from \textbf{\pa{ਧੰਨਵਾਦ}}. \textbf{\pa{ਦੁੱਧ}} is said exactly as spelled: \emph{dudh}, flat, no tone. Its breathy \textbf{\pa{ਧ}} never lost its breath — it is still a full consonant, the way \textbf{\pa{ਖ}} stayed a full consonant in \emph{khāṇā}.

\begin{cousinweb}
\textbf{dudh} is inherited from Sanskrit \textbf{dugdha}, \textquotedblleft{}that which has been milked\textquotedblright{} — grammatically a passive participle, built on the verb \textbf{duh-}, \textquotedblleft{}to milk.\textquotedblright{} Its Proto-Indo-European root meant something close to \textquotedblleft{}to produce, to be of use,\textquotedblright{} but the trail to any secure English word breaks there; this lesson claims none.

That gives you the chapter's first honest sort: \textbf{pāṇī} and \textbf{dudh} are both inherited straight from Sanskrit, with a known root behind each; \textbf{cāh} is a loan from Chinese. Three requests, two different kinds of history.
\end{cousinweb}

\subsection*{Your turn}
\begin{itemize}
\raggedright
  \item \emph{Say it:} \textbf{dudh, kirpā karke} — milk, please
  \item \emph{Run through:} all three requests so far — \textbf{pāṇī}, \textbf{cāh}, \textbf{dudh}, each with \emph{kirpā karke}
  \item \emph{Sort:} inherited (\textbf{pāṇī}, \textbf{dudh}) against loaned (\textbf{cāh})
  \item \emph{Name:} the mark that doubles a consonant, and one earlier word that used it
\end{itemize}

\begin{tcolorbox}[breakable,colback=teal!4,colframe=teal!35!black,title={Before you move on}]
Ask for milk politely. (\emph{Dudh, kirpā karke.}) What does \emph{dugdha} mean, literally? (\textbf{\textquotedblleft{}That which has been milked.\textquotedblright{}}) Does \emph{dudh} carry a tone? (\textbf{No} — its breathy \textbf{\pa{ਧ}} is still fully pronounced.)

Sources: \href{https://en.wiktionary.org/wiki/\%E0\%A8\%A6\%E0\%A9\%81\%E0\%A9\%B1\%E0\%A8\%A7}{Wiktionary: \pa{ਦੁੱਧ}}; \href{https://en.wiktionary.org/wiki/\%E0\%A4\%A6\%E0\%A5\%81\%E0\%A4\%97\%E0\%A5\%8D\%E0\%A4\%A7}{Wiktionary: Sanskrit dugdha}.
\end{tcolorbox}
\section[roṭī]{\pa{ਰੋਟੀ} — a request pattern, run four times}

\label{lesson:PA-C10-roti}



\begin{quote}
Three drinks, one pattern: \textbf{pāṇī}, \textbf{cāh}, \textbf{dudh}, each followed by \emph{kirpā karke}. The fourth request is food, not drink, and it closes the chapter on a letter this book has never properly taught.
\end{quote}

\subsection*{You'll want to know first — one word}
\begin{quote}
\textbf{\pa{ਰੋਟੀ}} — \emph{roṭī} — \textbf{bread}
\end{quote}

\begin{quote}
\textbf{\pa{ਰੋਟੀ}, \pa{ਕਿਰਪਾ} \pa{ਕਰਕੇ।}} — \emph{roṭī, kirpā karke.} — \textbf{Bread, please.}
\end{quote}

\subsection*{The letters in this word}
\textbf{\pa{ਰ}} \emph{r}, \textbf{\pa{ੋ}} the \emph{o} sign, then \textbf{\pa{ਟ}} — a retroflex \emph{ṭ}, tongue curled back and stopping the air rather than letting it through. Chapter 3 named the whole \textbf{\pa{ਟ}-\pa{ਠ}-\pa{ਡ}} series in passing, while teaching \textbf{\pa{ਠੀਕ}}'s \textbf{\pa{ਠ}}; this is the first time \textbf{\pa{ਟ}} itself has carried a word. Then \textbf{\pa{ੀ}} long \emph{ī}.

\begin{cousinweb}
\textbf{roṭī} is inherited from Sanskrit \textbf{roṭikā}, \textquotedblleft{}a kind of bread.\textquotedblright{} That much is certain. What \emph{roṭikā} itself comes from is not: the standard comparative dictionaries weigh an unattested root meaning \textquotedblleft{}to strike against\textquotedblright{} against another meaning \textquotedblleft{}to crush,\textquotedblright{} and neither wins.

That gives this chapter its own three-way sort: \textbf{pāṇī} and \textbf{dudh} are inherited straight from Sanskrit with a known root behind each; \textbf{cāh} is a loan, traceable to Mandarin by a known road; \textbf{roṭī} is inherited too, but its own root is where recorded history simply stops.
\end{cousinweb}

\subsection*{Your turn}
\begin{itemize}
\raggedright
  \item \emph{Say it:} \textbf{roṭī, kirpā karke} — bread, please
  \item \emph{Run through:} all four requests — \textbf{pāṇī}, \textbf{cāh}, \textbf{dudh}, \textbf{roṭī}, each with \emph{kirpā karke}
  \item \emph{Sort:} inherited with a known root (\textbf{pāṇī}, \textbf{dudh}), loaned (\textbf{cāh}), inherited with an unknown root (\textbf{roṭī})
  \item \emph{Name:} which earlier chapter first mentioned the \textbf{\pa{ਟ}-\pa{ਠ}-\pa{ਡ}} series
\end{itemize}

\begin{tcolorbox}[breakable,colback=teal!4,colframe=teal!35!black,title={Before you move on}]
Does \emph{roṭikā}'s own root have a settled answer? (\textbf{No} — the dictionaries disagree and leave it open.) Run all four requests from this chapter. (\emph{Pāṇī}, \emph{cāh}, \emph{dudh}, \emph{roṭī}, each with \emph{kirpā karke}.) Which one carries a tone, and what kind? (\textbf{Cāh} — high and level.)

Sources: \href{https://en.wiktionary.org/wiki/\%E0\%A8\%B0\%E0\%A9\%8B\%E0\%A8\%9F\%E0\%A9\%80}{Wiktionary: \pa{ਰੋਟੀ}}; \href{https://en.wiktionary.org/wiki/\%E0\%A4\%B0\%E0\%A5\%8B\%E0\%A4\%9F\%E0\%A4\%BF\%E0\%A4\%95\%E0\%A4\%BE}{Wiktionary: Sanskrit roṭikā}.
\end{tcolorbox}
